\begin{enumerate}
    \item Concentrations initiales~:
          \begin{itemize}
              \item $n_0(\ch{Fe^{2+}})=\bigl[\ch{Fe^{2+}}\bigr]_0\cdot V=
                    \num{1.0e-2}\times\num{100e-3}=\qty{1.0e-3}{\mol}$
              \item $n_0(\ch{Hg^{2+}})=[\ch{Hg^{2+}}]_0\cdot V=
                    \num{2.0e-2}\times\num{100e-3}=\qty{2.0e-3}{\mol}$
          \end{itemize}
          \begin{minipage}{\linewidth}
              \begin{table}[H]
              \centering
              \setlength{\arrayrulewidth}{0.8pt}
              \renewcommand{\arraystretch}{1.5}
              \begin{tabularx}{\textwidth}{|p{2.2cm}|
                  >{\centering\arraybackslash}p{2.6cm}|C|C|
                  >{\centering\arraybackslash}p{1.7cm}|
                  >{\centering\arraybackslash}p{1.7cm}|}
                  \hline \multicolumn{2}{|c|}{Équation de la réaction} &
                         \multicolumn{4}{c|}{%
                             \ce{\hspace{6mm}
                                2 Fe$^{2+}$_{(aq)} \hspace{6mm} + \hspace{6mm}
                                2 Hg$^{2+}$_{(aq)} \hspace{3mm} -> \hspace{3mm}
                                2 Fe$^{3+}$_{(aq)} \hspace{1mm} + \hspace{1mm}
                                Hg$_2^{2+}$_{(aq)}}} \\
                  \hline États & Avancement & \multicolumn{4}{c|}{Quantité de matière (en mol)} \\
                  \hline État initial & 0 & $\num{1.0e-3}$ & $\num{2.0e-3}$ &
                         $0$ & $0$ \\
                  \hline En cours & $x$ & $\num{1.0e-3}-2x$ &
                         $\num{2.0e-3}-2x$ & $2x$ & $x$ \\
                  \hline État final & $x_{\max}$ & $\num{1.0e-3}-2x_{\max}$ &
                         $\num{2.0e-3}-2x_{\max}$ & $2x_{\max}$ & $x_{\max}$ \\
                  \hline
              \end{tabularx}
          \end{table}
          \end{minipage}
          \textbf{(1pt)}
    \item 
          \begin{itemize}
              \item Si $\ch{Fe^{2+}}$ est le réactif limitant~:
                    $\num{1.0e-3}-2x_{\max}$
                    $\Rightarrow x_{\max}^{(1)}=\qty{5.0e-4}{\mol}$  
              \item Si $\ch{Hg^{2+}}$ est le réactif limitant~:
                    $\num{2.0e-3}-2x_{\max}$
                    $\Rightarrow x^{(2)}_{\max}=\qty{1.0e-3}{\mol}$  
          \end{itemize}
          $x^{(1)}_{\max}<x^{(2)}_{\max}$ $\Rightarrow$
          Le réactif limitant est \ch{Fe^{2+}}, donc~:
          $x_{\max}=\qty{5.0e-4}{\mol}$. \textbf{(0,5pt)}
        \item Le temps de demi-réaction est le temps nécessaire pour que
              l'avancement atteigne la moitié de sa valeur
              finale.~\textbf{(0,5pt)}

              Réaction totale ($x_{f}=x_{\max}$)~:
              $x(t_{1/2})=\frac{x_{\max}}{2}=\qty{2.5e-4}{\mol}$

              Graphiquement~: $t_{1/2}=\qty{14}{\h}$.~\textbf{(0,5pt)}
        \item La vitesse volumique est
              $v=\dfrac{1}{V}\cdot\dfrac{\mathrm{d}x}{\mathrm{d}t}$.~\textbf{(0,5pt)}

              À $t=0$,
              $\left.\frac{\mathrm{d}x}{\mathrm{d}t}\right|_{t=0}=
              \frac{(0,5-0)\times 10^{-3}}{20-0}=
              \qty{2.5e-5}{\mol\per\min}$.~\textbf{(0,25pt)}

              Donc~:
              $v_{t=0}=\frac{1}{V}
              \left.\frac{\mathrm{d}x}{\mathrm{d}t}\right|_{t=0}=
              \frac{1}{\qty{100e-3}{\L}}\times\qty{2.5e-5}{\mol\per\h}=
              \qty{2.5e-4}{\mol\per\L\per\h}$~\textbf{(0,25pt)}
        \item La vitesse volumique de réaction \textbf{diminue} au cours du
              temps, car les concentrations des réactifs diminuent. 
              La pente de la courbe $x(t)$ décroît jusqu'à devenir nulle lorsque
              $x$ atteint $x_{\max}$.
              \textbf{(0,5pt)}
    \end{enumerate}




